\chapter{Scheme-independent selection spin-ветви}

Полный base-\(K\) determinant не фиксирует абсолютный vacuum из-за
локальных finite counterterms. Однако разность determinant между двумя
flat/spin branches с одинаковыми локальными heat-kernel coefficients от
этой свободы не зависит. Поэтому дискретный branch selection является
следующим допустимым вопросом.

\section{Почему две ветви \texorpdfstring{\(\mathbb{RP}^3\)}{RP3} не различаются}

Для двух spin-структур круглого \(\mathbb{RP}^3\) signed Dirac spectra
обмениваются местами. При обозначениях \(\tau_\pm\) абсолютные уровни равны
\begin{equation}
 |\lambda_k|=\frac{k+3/2}{R_3},
 \qquad
 d_k=(k+1)(k+2),
 \qquad k\ge0,
\end{equation}
но знак уровня чередуется противоположно для \(\tau_+\) и \(\tau_-\).
Следовательно,
\begin{equation}
 \Spec D_{\mathbb{RP}^3,\tau_+}^2
 =\Spec D_{\mathbb{RP}^3,\tau_-}^2
\end{equation}
вместе с кратностями.

Vectorlike effective action двух физических Dirac-пар зависит от
\(D_K^2+\chi^2\). Поэтому его вещественный модуль не различает
\(\tau_+\) и \(\tau_-\). Их eta-инварианты
\begin{equation}
 \eta_D(\tau_\pm)=\pm\frac14
\end{equation}
также не снимают вырождение: determinant phases сокращаются между
vectorlike conjugate blocks.

\begin{proposition}[\(\mathbb{RP}^3\)-spin degeneracy]
Нетривиальный torsion character меняет signed branch, но не квадрат
спектра. В текущей vectorlike архитектуре две spin-структуры
\(\mathbb{RP}^3\) точно вырождены по вещественному one-loop action.
\end{proposition}

\section{Относительный determinant вдоль \texorpdfstring{\(S^1\)}{S1}}

Пусть окружная ветвь задаётся
\begin{equation}
 p_m=\frac{m+\beta}{R_1},
 \qquad
 \beta\in\mathbb R/\mathbb Z.
\end{equation}
Периодической и анти-периодической spin-структурам соответствуют
\(\beta=0\) и \(\beta=1/2\). Для одной трёхмерной Dirac shell положим
\begin{equation}
 \rho_k
 =\frac{R_1}{R_3}
 \sqrt{(k+3/2)^2+(\chi R_3)^2}.
\end{equation}
Точная разность окружных log-determinants равна
\begin{equation}
 \mathcal I_{\rho_k}(\beta)
 =\log
 \frac{\cosh(2\pi\rho_k)-\cos(2\pi\beta)}
 {\cosh(2\pi\rho_k)-1}.
\end{equation}
Локальная zero-winding часть не зависит от \(\beta\) и сокращается.

С учётом product-Clifford doubling и двух физических Dirac-пар получаем
\begin{equation}
 \Delta\Gamma_f(\beta)
 =-2\sum_{k=0}^{\infty}
 (k+1)(k+2)\,\mathcal I_{\rho_k}(\beta).
 \label{eq:fermion-spin-relative-action}
\end{equation}
Ряд абсолютно и экспоненциально сходится при всех конечных положительных
\(R_1/R_3\).

\section{Точный знак и выбор ветви}

Для любого \(\rho>0\)
\begin{equation}
 \mathcal I_\rho(\beta)\ge0,
\end{equation}
причём равенство достигается при \(\beta=0\pmod1\), а максимум --- при
\(\beta=1/2\pmod1\). Из фермионного минуса в
\eqref{eq:fermion-spin-relative-action} немедленно следует
\begin{equation}
 \Delta\Gamma_f(1/2)<\Delta\Gamma_f(0)=0.
\end{equation}
Более того, \(\beta=0\) является локальным максимумом, а
\(\beta=1/2\) --- глобальным минимумом relative fermion action.

При \(\chi R_3=1\) и \(R_1/R_3=1\) машинная сумма даёт
\begin{equation}
 \Delta\Gamma_f(1/2)
 =-0.0001948280\ldots.
\end{equation}
Знак сохраняется при всех проверенных \(\chi R_3\) и отношениях радиусов;
аналитическая положительность \(\mathcal I_\rho\) показывает, что это не
численный эффект cutoff.

\begin{theorem}[Условный anti-periodic selection]
Если functional integral суммирует две spin-структуры \(S^1\) с одинаковым
предварительным весом, fermion determinant текущего parent block строго
выбирает anti-periodic branch \(\beta=1/2\). Результат не зависит от
\(\mu\), \(\lambda_2\), \(\lambda_4\) и не использует наблюдаемые данные.
\end{theorem}

\begin{remark}[Граница утверждения]
Spin-структура может задаваться как внешнее superselection datum. В таком
чтении determinant только ранжирует две теории и не переводит одну в
другую динамически. Теорема становится физическим selection-механизмом
лишь при заранее объявленной сумме по spin structures или при dynamical
flat Wilson variable, содержащей точки \(\beta=0,1/2\).
\end{remark}

\section{Что закрыто и что остаётся}

Получено scheme-independent ранжирование дискретных ветвей без нового
коэффициента. При дополнительной гипотезе равновесной суммы оно даёт
\begin{equation}
 \boxed{\beta_{S^1}=\frac12.}
\end{equation}
Само ранжирование ещё не доказывает наличие spin-sum measure. Кроме того,
determinant modulus оставляет двойное вырождение
\(\tau_+\leftrightarrow\tau_-\) на \(\mathbb{RP}^3\), а абсолютный
\(\chi R\) по-прежнему зависит от renormalization conditions. Поэтому
следующий гейт обязан сначала вывести либо отвергнуть меру суммы по
spin-структурам. Только после этого разрешён поиск branch difference в
orientation/Pfaffian sector.

\begin{nogocriterion}[Запрет скрытой суммы по spin structures]
Нельзя объявлять anti-periodic branch динамически выбранной, если в
определении partition function spin-структура была заранее фиксирована.
Правило суммирования и относительные topological weights должны быть частью
parent theory до determinant-сравнения.
\end{nogocriterion}

Машинный аудит сохранён в
\path{s2t_v4_spin_structure_relative_determinant_gate_results.json}.